\documentclass[12pt, a4paper]{article}
% --- Encoding & language ---
\usepackage[utf8]{inputenc}
\usepackage[english]{babel}
% --- Math ---
\usepackage{amsmath, amsthm, amsfonts, amssymb}
\usepackage{mathtools}
\usepackage{bm}
% --- Layout & typography ---
\usepackage[margin=1in]{geometry}
\usepackage[expansion=false]{microtype}
\usepackage{parskip}
\usepackage{minted}
% --- References & links ---
\usepackage{hyperref}
\usepackage{footnotehyper}
\usepackage{cleveref}

\usepackage[round]{natbib}
\bibliographystyle{plainnat}

\usepackage{algorithm}
\usepackage{algpseudocode}

% --- Graphics---
\usepackage{graphicx}
\usepackage{tcolorbox}

% -----------------------------------------------
% Theorem environments
% -----------------------------------------------
\newtheoremstyle{spaced}
{1em}{1em}{\itshape}{}{\bfseries}{.}{0.5em}{}
\newtheoremstyle{spaceddef}
{1em}{1em}{\normalfont}{}{\bfseries}{.}{0.5em}{}
\theoremstyle{spaced}
\newtheorem{theorem}{Theorem}[section]
\newtheorem{lemma}[theorem]{Lemma}
\newtheorem{proposition}[theorem]{Proposition}
\newtheorem{corollary}[theorem]{Corollary}
\theoremstyle{spaceddef}
\newtheorem{definition}[theorem]{Definition}
\newtheorem{example}[theorem]{Example}
% Section numbers in equations
\numberwithin{equation}{section}

\newtheorem{remarkinner}[theorem]{Remark}
\newtheorem{statementinner}[theorem]{Statement}


\newenvironment{remark}
{\savenotes
    \begin{tcolorbox}[colback=white, colframe=black!80, arc=4pt, boxrule=0.5pt]%
        \begin{remarkinner}}
        {\end{remarkinner}\end{tcolorbox}\spewnotes}

\newenvironment{statement}
{\savenotes
    \begin{tcolorbox}[colback=white, colframe=black!80, arc=4pt, boxrule=0.5pt]%
        \begin{statementinner}}
        {\end{statementinner}\end{tcolorbox}\spewnotes}
% -----------------------------------------------
% Convenience macros
% -----------------------------------------------
\newcommand{\R}{\mathbb{R}}
\newcommand{\Z}{\mathbb{Z}}
\newcommand{\N}{\mathbb{N}}
% -----------------------------------------------
% Document metadata
% -----------------------------------------------
\title{Notes on Ensemble Smoothers}
\author{Tommy Odland}
\date{\today}
\setcounter{tocdepth}{1}
% -----------------------------------------------
% Document
% -----------------------------------------------
\begin{document}
\maketitle

\begin{abstract}
    \noindent
This article contains my notes on ensemble smoother algorithms.
They are used to approximate a posterior distribution using a set of samples.
By far the most common use-case is when the model is black-box and expensive to evaluate, for instance with petroleum reservoir models.

The notes are organized as follows: first we examine the Gauss-linear case, where we can compute an analytical posterior in arbitrarily high dimensions.
The equation that governs the Gaussian posterior is the foundation for ensemble methods.
Sadly reservoir models are not linear, so next we derive the Ensemble Smoother (ES)---an algorithm based on the Gauss-linear case.
Finally we examine EnIF, which is similar to ES but works with precision matrices rather than covariances, and leans on the statistical theory of high-dimensional estimation.
\end{abstract}

\tableofcontents

\section{Introduction}
\label{sec:introduction}

The Ensemble Kalman Filter was introduced by \citet{Evensen1994} and the Ensemble Smoother was introduced by \citet{VanLeeuwen1996}.
Four years later, in \citet{Burgers1998}, it was shown that the original update equation was wrong---the covariance did not match in the Gauss-linear case.\footnote{If you want to jump ahead, the issue was that in \cref{eqn:ensemble_smoother_update} the original formulation had $\bm{d} - \bm{y}$ in the innovation term instead of the correct $\bm{d} - (\bm{y} + \bm{\epsilon}_d)$}

New methods were proposed to deal with non-linearity.
One such variant is the
Iterative Ensemble Smoother (IES) of \citet{chen2012ensemble,chen2013levenberg}.
The idea behind IES is to move towards the maximum of the posterior in several small steps instead of a single large one, like a reduced step length optimization method.
Another variant is the Ensemble Smoother with Multiple Data Assimilation (ESMDA), proposed in \citet{emerick2013ensemble}.
It also makes several small steps, but shortens step length implicitly by increasing the covariance matrix of the observations.
Either way, the core idea is that linearizing, taking a small step, re-linearizing, taking another step, etc. is better than one linearization followed by one large step.

In high dimensions with few samples, empirical correlation matrices are mostly noise and very little signal.
All of the methods above suffer from this issue.
Many procedures have been proposed, such as setting all correlations that are below a $3 / \sqrt{r}$ to zero \citep{vossepoel2025adaptive}.
Another idea is to look towards the statistics of high-dimensional models, see for instance \citet[Chapter~18]{hastie2009elements}.
This is what \citet{lunde2025ensemble} does in the Ensemble Information Filter (EnIF), which uses graph properties to estimate sparse precision matrices and high-dimensional regression to estimate a sparse linearization of the forward model.

The papers above can be a bit tough to read:
notation can be a bit inconsistent, some proofs are not as detailed as I'd like, sometimes there is no code to look at, etc.
That is why I have decided to write these notes; maybe they can help clear up some confusion.
If nothing else, you'll get a second perspective.
You can find well-documented and production-ready Python code for many of these algorithms at the GitHub repository  \href{https://github.com/equinor/iterative_ensemble_smoother/}{github.com/equinor/iterative\_ensemble\_smoother/}.

\section{The Gaussian posterior}

% m = parameters, n = responses, r = realizations

Let $\bm{x} \in \R^m$ be a vector of \emph{parameters}.
The \emph{forward model} is denoted as $h: \R^m \to \R^n$, and applying it to the $m$ parameters returns a vector of $n$ \emph{responses} $\bm{y} \in \R^n$.

The \emph{observation} vector $\bm{d} \in \R^n$ is equal to the response plus a random noise component $\bm{\epsilon}_d$, which is often called the \emph{observation noise} or \emph{measurement noise}.
The full equation is then $\bm{d} = h(\bm{x}) + \bm{\epsilon}_d$,
where $\bm{\epsilon}_d \sim \mathcal{N}(\bm{0}, \Sigma_{\epsilon})$.

Think for instance of a petroleum reservoir model.
The parameters $\bm{x}$ are uncertain static properties of the subsurface, such as log-permeability and porosity in every cell of the simulation grid.
The forward model $h$ is a reservoir flow simulator.
It solves the fluid flow equations over the grid and reports predicted production data $\bm{y}$, for instance oil rate and bottom-hole pressure at each well over the producing history.
A single evaluation of $h$ may take hours of compute, and we only have access to its inputs and outputs.
The match is never exact: gauges drift, commingled flow must be allocated back to individual wells, and the grid only coarsely resembles the true geology, so the real-world observation $\bm{d}$ departs from the simulator response by a noise term $\bm{\epsilon}_d$.
The inverse problem of \emph{history matching} problem is: given the observed production $\bm{d}$, how can we update our knowledge of the parameters $\bm{x}$?
Characterizing the posterior distribution of $\bm{x}$ given $\bm{d}$ is the subject of the rest of these notes.

\subsection{Exact posterior distribution $p(\bm{x} \mid \bm{d})$}

In this section we assume that the prior distribution over parameters $p(\bm{x})$ is multivariate Gaussian, and that the forward model $h$ is linear.
Then $\bm{d} = H \bm{x} + \bm{\epsilon}_d = \bm{y} + \bm{\epsilon}_d$, where $H \in \R^{n \times m}$ is a matrix.
The posterior distribution $p(\bm{x} \mid \bm{d})$ is also multivariate Gaussian.

In words, we will answer ``Suppose I have a prior belief over parameters $p(\bm{x})$ and I have observed an observation $\bm{d}$---how can I update my prior to a posterior?''

\begin{remark}[Response covariance]
    \label{remark:response_covariance}
    Recall from \citet[Section~6.2]{petersen2012matrix} that $\operatorname{Cov}\left[A\bm{x}, B\bm{y}\right] = A \operatorname{Cov}\left[\bm{x}, \bm{y}\right] B^T$ when $\bm{x}$ and $\bm{y}$ are random variables.
    When both arguments are identical, we'll use a shorthand notation:  $\operatorname{Cov}\left[\bm{x}\right] := \operatorname{Cov}\left[\bm{x}, \bm{x}\right] = \operatorname{Var}\left[\bm{x}\right]$.
    Similarly, when there is no ambiguity we'll simply write $\Sigma_x$ instead of $\Sigma_{xx}$ for the covariance of $\bm{x}$.

    The covariance in $\bm{d}$ is the covariance of $\bm{x}$, propagated through the linear map $H$, plus the covariance of the observation noise $\bm{\epsilon}_d$:
    \begin{align*}
        \Sigma_{d} &=
        \operatorname{Cov}\left[ \bm{d} \right]
        = \operatorname{Cov}\left[ H\bm{x} + \bm{\epsilon}_d \right] \\
        &= \operatorname{Cov}\left[ H\bm{x}  \right]
        + \operatorname{Cov}\left[ \bm{\epsilon}_d \right]
        = H \operatorname{Cov}\left[\bm{x}, \bm{x}\right] H^T
        + \Sigma_{\epsilon} \\
        & = H \Sigma_x H^T
        + \Sigma_{\epsilon}.
    \end{align*}
    Similarly, the cross-covariance can be expressed as:
\begin{align*}
    \Sigma_{xd} &=
    \operatorname{Cov}\left[ \bm{x}, \bm{d} \right] =
    \operatorname{Cov}\left[ \bm{x}, H\bm{x} + \bm{\epsilon}_d \right]  \\
    &=
    \operatorname{Cov}\left[ \bm{x}, \bm{x} \right] H^T = \Sigma_x H^T
\end{align*}
In both equations above we used the fact that $\bm{\epsilon}_d$ is independent of both $\bm{x}$ and $\bm{y}$,
which also means that $\Sigma_{xd} = \Sigma_{xy}$.
\end{remark}

Suppose the prior over parameters $\bm{x}$ and the observation $\bm{d}$, conditioned on the parameters, are given by multivariate Gaussians:
\begin{align*}
    p(\bm{x}) &= \mathcal{N}\left( \bm{x} \mid \bm{\mu}_x, \Sigma_x \right) \\
    p(\bm{d} \mid \bm{x}) &= \mathcal{N}( \bm{d} \mid H\bm{x}, \Sigma_{\epsilon})
\end{align*}
Then the posterior distribution over parameters, conditioned on the response, is also a multivariate Gaussian $p(\bm{x} \mid \bm{d}) = \mathcal{N} (\bm{\mu}_{x \mid d}, \Sigma_{x \mid d})$.
The posterior mean and covariance are:
\begin{align}
    \bm{\mu}_{x \mid d} &=
    \Sigma_{x \mid d} \left(
    H^T \Sigma_\epsilon^{-1} \bm{d} + \Sigma_x^{-1} \bm{\mu}_x
    \right)
    \label{eqn:posterior_mean_precision}
    \\
    \Sigma_{x \mid d} &=
    \left( \Sigma_x^{-1} + H^T \Sigma_\epsilon^{-1} H  \right)^{-1}
        \label{eqn:posterior_cov_precision}
\end{align}
The equations above are well known, and can for instance be found in \citet[Section~2.3.3]{BISH11}.\footnote{We can derive the posterior mean $\bm{\mu}_{x \mid d}$ by maximizing the log-posterior $\log ( p(\bm{x} \mid \bm{d}))$, which is proportional to  $\log ( p(\bm{d} \mid \bm{x}) \, p(\bm{x}))$ by Bayes' theorem.
Simply take the gradient of
$(\bm{d} - H \bm{x})^T \Sigma_\epsilon^{-1} (\bm{d} - H \bm{x})
+
(\bm{x} - \bm{\mu}_x)^T \Sigma_x^{-1} (\bm{x} - \bm{\mu}_x)
$ with respect to $\bm{x}$, equate it to $\bm{0}$ and solve.}
In the equations above, the posterior mean $\bm{\mu}_{x \mid d}$ is expressed in terms of precision matrices rather than covariance matrices.
% Conditional Gaussians are more naturally expressed in terms of precision, while marginals are more naturally expressed in terms of covariance.

\begin{remark}[Posterior covariance]
    \label{remark:posterior_covariance}
    We show the Schur complement form of the posterior covariance $\Sigma_{x \mid d}$.
    Starting with \cref{eqn:posterior_cov_precision}, we first use the Woodbury identity, then the results from \cref{remark:response_covariance}:
\begin{align*}
    \Sigma_{x \mid d} &=
    \left( \Sigma_x^{-1} + H^T \Sigma_\epsilon^{-1} H  \right)^{-1} \\
    &= \Sigma_x - \Sigma_x H^T
    \left(
    \Sigma_\epsilon + H \Sigma_x H^T
    \right)^{-1}
    H \Sigma_x \\
    &= \Sigma_x - \Sigma_{xd}
    \Sigma_d^{-1}
    \Sigma_{dx} \\
\end{align*}
\end{remark}

In \cref{remark:posterior_covariance} we derived the Schur complement form of the update equation.
We will now derive the Kalman gain form, which is given by:
\begin{align}
    \nonumber
    \bm{\mu}_{x \mid d} &= \bm{\mu}_{x} + K (\bm{d} - H \bm{\mu}_{x}) \\
    \label{eqn:kalman_gain}
    K &:= \Sigma_x H^T
    \left(
    \Sigma_\epsilon + H \Sigma_x H^T
    \right)^{-1} = \Sigma_{xd} \Sigma_d^{-1}
\end{align}
In the equation above, $H \bm{\mu}_{x}$ is the predicted mean, $(\bm{d} - H \bm{\mu}_{x})$ is the \emph{innovation} term and $K$ is the Kalman gain.
In words: the posterior equals the prior, plus a linear combination of the innovation (residuals).

There are many ways to derive the Kalman update equation.
One way is to look at \citet[Section~2.3.1]{BISH11}, but below we will take a different route.

The first step is to show that $K = \Sigma_{x \mid d}H^T \Sigma_{\epsilon}^{-1}$.
To see this, we will start by assuming that both expression are equal:
\begin{equation}
    \label{eqn:push_through_kalman}
    \Sigma_{x \mid d}H^T \Sigma_{\epsilon}^{-1} =
    K =
    \Sigma_x H^T
    \left(
    \Sigma_\epsilon + H \Sigma_x H^T
    \right)^{-1}
\end{equation}
We right-multiply both sides with $\left(
\Sigma_\epsilon + H \Sigma_x H^T
\right)$, then left-multiply both sides with $\Sigma_{x \mid d}^{-1}$ and finally expand both sides to show that they both equal $H^T + H^T \Sigma_\epsilon^{-1} H \Sigma_x H^T$.\footnote{A more direct, constructive way of showing this is to appeal to the \emph{push-through identity}, see for instance the Wikipedia page on the Woodbury matrix identity.}

The second step is to realize that
\begin{align*}
    \Sigma_{x \mid d}\Sigma_{x}^{-1} &=
    \Sigma_{x \mid d}
    \left( \Sigma_{x \mid d}^{-1} - H^T \Sigma_{\epsilon}^{-1} H \right) \\
    &= I - \Sigma_{x \mid d}  H^T \Sigma_{\epsilon}^{-1} H =  I - KH.
\end{align*}

With these two results in hand, the update equation for the mean becomes
\begin{align*}
    \bm{\mu}_{x \mid d} &=
    \Sigma_{x \mid d} \left(
    H^T \Sigma_\epsilon^{-1} \bm{d} + \Sigma_x^{-1} \bm{\mu}_x
    \right) \\
    &=
    \Sigma_{x \mid d}
    H^T \Sigma_\epsilon^{-1} \bm{d} +  \Sigma_{x \mid d} \Sigma_x^{-1} \bm{\mu}_x
     \\
    & =K \bm{d} + (I - KH) \bm{\mu}_x =  \bm{\mu}_x + K (\bm{d} - H\bm{\mu}_x )
\end{align*}

In summary then, the Kalman-style update equations are:
\begin{align*}
    \bm{\mu}_{x \mid d}
    & = \bm{\mu}_x + K (\bm{d} - H\bm{\mu}_x )
    \\
    \Sigma_{x \mid d} &= (I - KH) \Sigma_x
\end{align*}
We have seen that there are different way to express the update equations: the precision form of \cref{eqn:posterior_cov_precision}, the Schur form of \cref{remark:posterior_covariance} and the Kalman form above.
Which one should we use?
It depends on what we wish to show or prove, or how we want to structure our computations.
For instance, the Ensemble Smoother---which we will discuss next---uses a Kalman style update with covariances.

\section{The Ensemble Smoother}
% Above we described the posterior as a multivariate Gaussian specified by the mean $\bm{\mu}_{x \mid d}$ and covariance $\Sigma_{x \mid d}$.
% What if we instead describe the posterior by a set of samples?
The Ensemble Smoother (ES) describes the posterior with an \emph{ensemble} of \emph{realizations} (i.e., a set of samples) instead of $\bm{\mu}_{x \mid d}$ and $\Sigma_{x \mid d}$.
First we draw $r$ realizations $\{\bm{x}_j\}_{j=1}^{r}$ from the prior distribution $p(\bm{x})$.
Then we propagate them through the forward model $h$ and obtain $\bm{y}_j = h(\bm{x}_j)$ for each realization $j$.
Finally, we perform \emph{data assimilation}; a Kalman-style update that transforms the prior ensemble to the posterior ensemble.

The upside of ES is that we don't need knowledge of the forward model $h$ to compute a posterior; $h$ can be a black-box function.
The downside is that the update equation is derived assuming (1) a linear $h$, (2) Gaussian distributions and (3) infinitely many realizations.
ES won't actually sample the true, analytical posterior except in the happy case when conditions (1), (2) and (3) are all satisfied.
Instead we approximate it.\footnote{This approximation comes without any guarantees, so engineering judgment is often required. We potentially perform an update so full of noise that keeping the prior would've been better.}
Since forward model is expensive to compute, the number of realizations $r$ is much smaller than the number of parameters $m$ and responses $n$.
This leads to poor estimates of covariances, a phenomenon often called \emph{spurious correlations}.

We will now derive the update equation for ES, which computes the posterior realizations given the prior realizations.
First, some notation: let $X \in \R^{m \times r}$ be a matrix with $r$ realizations (ensemble members); one per column.\footnote{In statistics and machine learning, samples are typically associated with rows and features with columns. Here that convention is transposed; each parameter gets a row and each realization a column.}
Let $Y = h(X) = [ h(\bm{x}_1) \mid h(\bm{x}_2) \mid \cdots \mid h(\bm{x}_r)] \in \R^{n \times r}$ be a matrix of responses, where $h$ is applied to each column.

\begin{statement}
The matrix form of the ES update equation, which updates every realization (column) in $X$, is
\begin{equation}
    \label{eqn:ensemble_smoother_update}
    X_{\text{posterior}}
    = X_{\text{prior}} +
    \hat{\Sigma}_{xy}
    (
    \Sigma_\epsilon + \hat{\Sigma}_y
    )^{-1} (\bm{d} \bm{1}^T - (Y + E))
\end{equation}
where each column of $E$ is independently sampled observation $\bm{\epsilon}_d$.
In other words, $E_{\cdot, j} \sim \mathcal{N}(\bm{0}, \Sigma_\epsilon)$ for every column $j$.
The matrices $\hat{\Sigma}_{xy}$ and $\hat{\Sigma}_y$ are estimates of the cross-correlation and the correlation.
The most common estimate for, e.g.,  $\hat{\Sigma}_{xy}$ is the empirical cross correlation matrix $\bar{X}\bar{Y}^T / (r-1) \in \R^{m \times n}$,
where $\bar{X}$ is a centered $X$ in which every parameter (row) has its mean subtracted.\footnote{Some authors write the centered matrix $\bar{X}$ as $X (I - \bm{1} \bm{1}^T / r)$.
While centering can be expressed as matrix multiplication, it is both hard to read and an unwise implementation strategy.
}
$\bar{Y}$ is a centered $Y$ and $\hat{\Sigma}_{y}$ is computed similarly.
\end{statement}
To derive \cref{eqn:ensemble_smoother_update}, we start by examining the Gaussian update equation for the mean:
\begin{align*}
    \bm{\mu}_{x \mid d}
    & = \bm{\mu}_x + K (\bm{d} - H\bm{\mu}_x ) \\
    & = \bm{\mu}_x + \Sigma_x H^T
    \left(
    \Sigma_\epsilon + H \Sigma_x H^T
    \right)^{-1} (\bm{d} - H\bm{\mu}_x )
\end{align*}
% Since we can no longer assume knowledge of $H$,
We use the results from \cref{remark:response_covariance} and substitute
$\Sigma_x H^T = \Sigma_{xy}$
and
$H \Sigma_x H^T = \operatorname{Cov}\left[ H\bm{x}  \right]
= \operatorname{Cov}\left[ \bm{y}  \right]
= \Sigma_y$.
This gives us Kalman gain in the form $K = \Sigma_{xy}
\left(
\Sigma_\epsilon + \Sigma_y
\right)^{-1}$.

The update equation for each realization $j$, is given by
\begin{equation*}
    \bm{x}^{j}_{\text{posterior}}
    =
    \bm{x}^{j}_{\text{prior}} + K(\bm{d} - (\bm{y}^j + \bm{\epsilon}_d^j)),
\end{equation*}
where $\bm{\epsilon}_d^j \sim \mathcal{N}(\bm{0}, \Sigma_\epsilon)$.
To see why this is the correct update and why we have to perturb the response $\bm{y}^j$ with observation noise $\bm{\epsilon}_d^j$, we have to prove that in the Gauss-linear case, this update equation recovers the same posterior mean and posterior covariance as the analytical equations in \cref{eqn:posterior_mean_precision,eqn:posterior_cov_precision}.
To this end, suppose that $h$ is linear, so $\bm{y}^j = H \bm{x}^{j}_{\text{prior}}$.

We update each realization with the equation
\begin{align*}
    \bm{x}_{\text{posterior}}
    &=
    \bm{x}_{\text{prior}} + K(\bm{d} - (H \bm{x}_{\text{prior}} + \bm{\epsilon}_d)) \\
    \bm{x}_{\text{prior}} &\sim \mathcal{N}(\bm{\mu}_x, \Sigma_x) \\
    \bm{\epsilon}_d &\sim \mathcal{N}(\bm{0}, \Sigma_\epsilon)
\end{align*}

First we show that the posterior mean equals \cref{eqn:posterior_mean_precision}:
\begin{align*}
    \mathbb{E} \left[ \bm{x}_{\text{posterior}}  \right]
    &= \mathbb{E} \left[ \bm{x}_{\text{prior}} +
    K(\bm{d} - (H \bm{x}_{\text{prior}} + \bm{\epsilon}_d))  \right] \\
    &= \mathbb{E} \left[ \bm{x}_{\text{prior}}  \right]
    +
    K
    \left(
    \mathbb{E} \left[ \bm{d}  \right]
    -
    \mathbb{E} \left[  H \bm{x}_{\text{prior}} \right]
    -
    \mathbb{E} \left[  \bm{\epsilon}_d  \right]
    \right) \\
    &=
    \bm{\mu}_x + K \left(
     \bm{d}
    -
    H \mathbb{E} \left[  \bm{x}_{\text{prior}}  \right]
    -
    \bm{0}
    \right) \\
    &=
    \bm{\mu}_x + K(\bm{d}  - H \bm{\mu}_x )
\end{align*}

Second we show that the posterior covariance equals \cref{eqn:posterior_cov_precision}.
More specifically, it will be equal to the Schur complement form derived in \cref{remark:posterior_covariance}.
Below we use the fact that $K\bm{d}$ is constant and that that $\operatorname{Cov} \left[ \bm{x}_{\text{prior}},  \bm{\epsilon}_d \right] = 0$ when we distribute the covariance over the terms.
\begin{align*}
    \operatorname{Cov} \left[ \bm{x}_{\text{posterior}}  \right]
    &= \operatorname{Cov} \left[ \bm{x}_{\text{prior}} +
    K(
    \bm{d} - (H \bm{x}_{\text{prior}} + \bm{\epsilon}_d)
    )  \right] \\
    &= \operatorname{Cov} \left[
    (I - KH) \bm{x}_{\text{prior}}
    + K \bm{d}
    - K \bm{\epsilon}_d
     \right] \\
     &= \operatorname{Cov} \left[
     (I - KH) \bm{x}_{\text{prior}}
     \right] +
     \operatorname{Cov} \left[
     -K \bm{\epsilon}_d
     \right]
     \\
     % this is called the Joseph form
     &= (I - KH) \Sigma_{x} (I - KH)^T + K \Sigma_{\epsilon} K^T \\
     &= \Sigma_{x} - \Sigma_{x} H^T K^T - K H \Sigma_{x} + K H \Sigma_{x} H^T K^T + K \Sigma_{\epsilon} K^T \\
     &= \Sigma_{x} -
     \underbrace{\Sigma_{x} H^T}_{\Sigma_{xd}} K^T -
     K \underbrace{H \Sigma_{x}}_{\Sigma_{dx}} +
     K \underbrace{\left( H \Sigma_{x} H^T + \Sigma_{\epsilon} \right) }_{\Sigma_{d}} K^T \\
     &= \Sigma_{x} - \Sigma_{xd} K^T - K \Sigma_{dx} +
     K \Sigma_{d} K^T
\end{align*}
The equalities shown in underbraces were all derived in \cref{remark:response_covariance}.
From \cref{eqn:kalman_gain}, we observe that $K \Sigma_{d} = (\Sigma_{xd} \Sigma_{d}^{-1}) \Sigma_{d} = \Sigma_{xd}$.
The last term, $K \Sigma_{d} K^T$, can therefore be re-written as $\Sigma_{xd} K^T$ and cancels the second term.
We are left with
\begin{equation*}
    \operatorname{Cov} \left[ \bm{x}_{\text{posterior}}  \right]
    = \Sigma_{x}  - K \Sigma_{dx} = \Sigma_{x} - \Sigma_{xd} \Sigma_{d}^{-1} \Sigma_{dx}
    = \Sigma_{x \mid d}.
\end{equation*}
We have shown that if we use \cref{eqn:ensemble_smoother_update} to update the ensemble, then the resulting posterior mean and covariance over the realizations matches the analytical posterior in the Gauss-linear case.
% Note again that this only holds when (1) $h$ is linear, (2) the prior and observation noise are Gaussian distributions and (3) we have infinitely many realizations.
% If any of these assumptions are violated, then the posterior is an approximation.

\subsection{Computation}

TODO: Here I could write about practical computation, since it differs a lot from the naive equation. but it is already well documented in the code.
% ES-MDA stands for Ensemble Smoother with Multiple Data Assimilation.
% In our framework, it can be derived as follows:

\subsection{Extensions}

TODO: Here i could write a very short section about extension: localization, etc


\section{The Ensemble Information Filter}

One issue with ES is that in \cref{eqn:ensemble_smoother_update} we have to estimate $\Sigma_{xy}$ and $\Sigma_{y}$.
When $m$ and $n$ are larger than $r$, empirical covariance is a poor estimator that picks up spurious correlations.

What if we avoid covariances and instead work with precision matrices and $H$?
After all, there is a substantial literature on high-dimensional regression that we can lean on to estimate $H$.
Also, parameters with conditional independence structure are often better described by precision matrices $\Lambda_x$ rather than covariance matrices $\Sigma_{x}$, see \citet[Chapter~2]{rue2005gaussian}.
The following ideas are due to \citet{lunde2025ensemble}:
(1) estimate $H$ with sparse high dimensional regression and
(2) think in terms of precision rather than covariance.
We'll refer to his work as the Ensemble Information Filter (EnIF).
The \emph{information} terminology means that we use precision rather than covariance.
We're discussing smoothers, so the \emph{filter} does not exactly match this document---but we'll stick with EnIF rather than invent a new abbreviation.

Let us approximate $\bm{y} = h(\bm{x})$ as $\bm{y} = H \bm{x} + \bm{r}$.
Here $\bm{r}$ is the residual (the non-linear part of $h$) and $H$ is a sparse matrix.\footnote{You might be wondering why we do not include a bias term in the approximation? It's not needed, since it will cancel in the innovation term in ES-style updates like \eqref{eqn:ensemble_smoother_update} anyway.}
% We do assume that the forward model $h$ is deterministic; running it on the identical inputs (parameters) produces identical outputs (responses).
% In that sense $\bm{r}$ is not a random variable, but we will treat it as random because it allows us to capture and model the quality of $H$ and adjust updates accordingly.
Assuming $\bm{r} \sim \mathcal{N}(\bm{0}, \Sigma_{r})$, then if $\Sigma_{r}$ is large the linear model is poor.
\footnote{The reason we did not include residuals $\bm{r}$ in the ES derivation is that they become part of $\Sigma_{y}$ in \eqref{eqn:ensemble_smoother_update} anyway.
Under the assumption $\bm{y} = H \bm{x}$ we have
$\Sigma_{y}
= \operatorname{Cov}\left[ \bm{y} \right]
= \operatorname{Cov}\left[ H \bm{x} \right]
= H \Sigma_{x} H^T$.
Had we instead assumed $\bm{y} = H \bm{x} + \bm{r}$, then
$\Sigma_{y}
= \operatorname{Cov}\left[ \bm{y} \right]
= \operatorname{Cov}\left[ H \bm{x} + \bm{r} \right]
= H \Sigma_{x} H^T + \Sigma_{r}$.
Either way, in the ES update equation \eqref{eqn:ensemble_smoother_update} we rely on an empirical $\Sigma_{y}$, which captures the residuals.
It is only now that we wish to rely on $H$ instead of $\Sigma_{y}$ that we have to decompose $\Sigma_{y}$ into variance due to $\bm{x}$ mapped through $H$ (the $H \Sigma_{x} H^T$ term) and variance due to residuals (the $\Sigma_{r}$ term).}

\begin{statement}
    The matrix form of the EnIF update equation, which updates every realization (column) in $X$, is
    \begin{align}
        \label{eqn:ensemble_smoother_update_inf}
        X_{\text{posterior}}
        &= X_{\text{prior}} +
        ( \Lambda_x + H^T \tilde{\Lambda} H )^{-1} H^T \tilde{\Lambda}
         (\bm{d} \bm{1}^T - (Y + E)) \\
         \label{eqn:ensemble_smoother_update_precision_def}
         \tilde{\Lambda} &= (\Sigma_{d \mid x})^{-1} =
         \left( \Sigma_{\epsilon} + \Sigma_{r} \right)^{-1}
    \end{align}
where $Y = h(X_{\text{prior}})$ applied column-wise and each column of $E$ is independently sampled observation $\bm{\epsilon}_d$.
In other words, $E_{\cdot, j} \sim \mathcal{N}(\bm{0}, \Sigma_\epsilon)$ for every column $j$.

$H$, $\Lambda_x$ and $\Sigma_{r}$ must be estimated.
The forward map $H$ is estimated by regressing $Y = h(X_{\text{prior}})$ on $X_{\text{prior}}$, typically with a sparse high dimensional method.
The precision $\Lambda_x$ must be estimated, taking into account the parameter ensemble $X$ as well as neighborhood (graph) information (sparsity pattern in $\Lambda_x$).
The residual covariance $\Sigma_{r}$ can be estimated as
$\operatorname{diag}\left( \operatorname{Cov}\left[ Y - HX \right] \right)$, where $H$ is the estimated forward map.
Correcting for degrees of freedom is an option when estimating $\Sigma_{r}$.
\end{statement}

We will now sketch the proof of \cref{eqn:ensemble_smoother_update_inf} being the correct update equation, in the sense that posterior mean and covariance matches under the Gauss-linear model.
The analytical model given the assumption $\bm{y} = H \bm{x} + \bm{r}$ becomes:
\begin{align*}
    p(\bm{x}) &= \mathcal{N}\left( \bm{x} \mid \bm{\mu}_x, \Sigma_x \right) \\
    p(\bm{d} \mid \bm{x}) &= \mathcal{N}( \bm{d} \mid H\bm{x}, \Sigma_{\epsilon} + \Sigma_{r})
\end{align*}
The posterior distribution $p(\bm{x} \mid \bm{d}) = \mathcal{N} (\bm{\mu}_{x \mid d}, \Sigma_{x \mid d})$ over parameters is given by
\begin{align}
    \bm{\mu}_{x \mid d} &=
    \Sigma_{x \mid d} \left(
    H^T \left( \Sigma_{\epsilon} + \Sigma_{r}  \right)^{-1} \bm{d} + \Sigma_x^{-1} \bm{\mu}_x
    \right)
    \label{eqn:posterior_mean_precision_inf}
    \\
    \Sigma_{x \mid d} &=
    \left( \Sigma_x^{-1} + H^T \left( \Sigma_{\epsilon} + \Sigma_{r}  \right)^{-1} H  \right)^{-1}
    \label{eqn:posterior_cov_precision_inf}
\end{align}
From \cref{eqn:push_through_kalman} and the result in \cref{remark:response_covariance}, we observe that we can write the Kalman gain as a function of precision matrices and $H$ as
\begin{equation*}
    K
    = \Sigma_{x \mid d} H^T \Sigma_{d \mid x}^{-1}
    = ( \Lambda_x + H^T \tilde{\Lambda} H )^{-1} H^T \tilde{\Lambda}
\end{equation*}
where $\tilde{\Lambda}$ is defined as in \cref{eqn:ensemble_smoother_update_precision_def}.
At this point it's quite easy to mimic the proof for the ES update equation, to show that \cref{eqn:ensemble_smoother_update_inf} also is correct, in the sense that
\begin{equation*}
    \mathbb{E} \left[ \bm{x}_{\text{posterior}}  \right] = \bm{\mu}_{x \mid d}
    \qquad
    \operatorname{Cov} \left[ \bm{x}_{\text{posterior}}  \right] = \Sigma_{x \mid d}
\end{equation*}
The two differences between EnIF and ES, under the gauss-Linear assumption, are that (1) the Kalman gain relies on precision and $H$ instead of covariance and (2) the covariance of $\bm{r}$ is accounted for.

\subsection{Computation}

TODO: This section is still a sketch...

Step 1:
Given $X_\text{prior}$ and possibly a sparsity pattern $G$, estimate the (possibly sparse) precision matrix $\Lambda_x$.

Step 2:
Compute $Y = h(X_\text{prior})$.
Then find a sparse $H$ so that $Y \sim HX$.
Compute the residual matrix $R = Y - HX$ and form it estimate
$\Sigma_r$, for instance by taking the variance over each row (response-residual) in $R$.

Step 3:
Solve
\begin{equation*}
( \Lambda_x + H^T \tilde{\Lambda} H ) X = H^T \tilde{\Lambda}
(D - Y )
\end{equation*}
or, equivalently, the column version for each column $j$
\begin{equation*}
    ( \Lambda_x + H^T \tilde{\Lambda} H ) \bm{x}_j = H^T \tilde{\Lambda}
    ((\bm{d} + \bm{\epsilon}_d) - \bm{y}_j )
\end{equation*}
If $\tilde{\Lambda}$ is diagonal, then exploit that fact.
Since $( \Lambda_x + H^T \tilde{\Lambda} H )$ is sparse, large and symmetric positive definite,
it can be solved with Cholesky or with Conjugate gradients.
% TODO: which one is better? any other good alternatives?
% if i do cholesky i can re-use the factorization for each column x_j.
% can i do this with conjugate gradients?

Step 4:

Set $X_{\text{posterior}}
= X_{\text{prior}} + X$

\begin{statement}
    content...
\end{statement}


\subsection{Extensions}

TODO: here i could write about estimating H and prec




% BIBLIOGRAPHY
\bibliography{references}
\end{document}
